\documentclass[12pt,a4paper]{article}

\usepackage[spanish,es-nodecimaldot]{babel}
\usepackage[utf8]{inputenc}
\usepackage[T1]{fontenc}
\usepackage{lmodern}
\usepackage{amsmath,amssymb}
\usepackage{graphicx}
\usepackage{booktabs}
\usepackage{hyperref}
\usepackage{enumitem}
\usepackage{geometry}
\usepackage{listings}
\usepackage{xcolor}

\geometry{margin=2.5cm}

\lstset{
  basicstyle=\footnotesize\ttfamily,
  breaklines=true,
  frame=single,
}

\begin{document}

\section{Investigación profunda para defensa (paquete multiarchivo)}

Este paquete está pensado para preparar una defensa completa del proyecto \textbf{TFG\_CYBER\_AI}
ante un tribunal y también para explicarlo a una persona sin base previa.

\subsection{Objetivo de este paquete}

\begin{itemize}
  \item cubrir de forma rigurosa los temas clave técnicos y metodológicos
  \item explicar cada parte con lenguaje accesible
  \item separar claramente lo \textbf{implementado}, lo \textbf{histórico} y lo \textbf{pendiente}
  \item evitar afirmaciones sin respaldo en código o artefactos de \verb|runs/|
\end{itemize}

\subsection{Relación con investigaciones previas}

Ya existen dos investigaciones profundas:

\begin{enumerate}
  \item \verb|docs/Personal Research/data-structure-and-canonical-schema-research-report.md|
  \item \verb|docs/Personal Research/qrdqn-research-report.md|
\end{enumerate}

Este nuevo paquete las \textbf{complementa} y las integra en una narrativa de defensa completa.

\subsection{Orden recomendado de lectura}

\begin{enumerate}
  \item \texttt{01-fundamentos-y-objetivo}
  \item \texttt{02-datos-esquema-canonico-y-preprocesado}
  \item \texttt{03-entorno-rl-algoritmo-y-entrenamiento}
  \item \texttt{04-validacion-y-lectura-de-resultados}
  \item \texttt{05-phase2-laboratorio-inferencia-y-riesgos}
  \item \texttt{06-glosario-y-preguntas-tribunal}
\end{enumerate}

\subsection{Mapa rápido de entradas técnicas del repositorio}

\begin{itemize}
  \item Esquema canónico: \verb|src/canonical_schema.py|
  \item Carga CICIDS2017: \verb|src/load_cicids2017.py|
  \item Entorno RL: \verb|src/rl_defender_env.py|
  \item Entrenamiento QRDQN: \verb|src/train_rl_defender.py|
  \item Validación A/B/C: \verb|src/validate_checks.py|
  \item Leave-one-exact-CSV-out: \verb|src/validate_leave_one_csv_out.py|
  \item Inferencia Phase~2 (mantenida): \verb|scripts/predict_real_traffic_v2.py|
  \item Utilidades de clipping: \verb|src/scaling_utils.py|
\end{itemize}

\end{document}
